% PDF 1.4 without object streams, for older viewers.
\pdfminorversion=4
\pdfobjcompresslevel=0
\documentclass[11pt]{article}
\usepackage[margin=2.3cm]{geometry}
\usepackage{amsmath,amssymb}
\usepackage{booktabs,array,longtable}
\usepackage[T1]{fontenc}
\usepackage{lmodern}
\usepackage{microtype}
\usepackage[hidelinks]{hyperref}

\newcommand{\Sq}{\mathrm{Sq}}
\newcommand{\F}{\mathbb{F}_2}
\newcommand{\A}{\mathcal{A}}
\DeclareMathOperator{\LT}{LT}
\DeclareMathOperator{\Ext}{Ext}
\setlength{\parindent}{0pt}
\setlength{\parskip}{0.4em}

\title{A small minimal resolution over the Steenrod algebra,\\ computed by Gröbner bases}
\author{Trace of the \texttt{groebner} crate (branch \texttt{groebner\_lin})}
\date{}

\begin{document}
\maketitle

\begin{abstract}
We follow, one reduction at a time, the computation of the minimal free resolution
$\dots \to C_3 \to C_2 \to C_1 \to C_0 \to \F$ of the ground field over the mod $2$ Steenrod
algebra $\A$ by the \texttt{groebner} crate, a port of Weinan Lin's Gröbner basis method. The
range is homological degree $s \le 3$ and stem $t - s \le 4$, so internal degrees $t \le 7$.
Every formula below is taken from the program's own data: the critical pairs it formed, the
elements it reduced, and the bases and differentials it produced. Only the bookkeeping of the
reduction steps was replayed outside the program, and the replay was checked to reproduce the
program's results exactly.
\end{abstract}

\section{Setting}

\paragraph{The algebra.}
We work in the Milnor basis $\Sq(R) = \Sq(r_1, r_2, \dots)$ of $\A$, writing $\Sq^r = \Sq(r)$.
Write $r_n = \sum_i d_{i,n} 2^i$ in binary. The digit $d_{i,n}$ records the power
$\xi_n^{2^i}$ of the dual basis, of degree $2^i(2^n - 1)$, and we record $\Sq(R)$ by the
monomial
\[ x^R = \prod_{d_{i,n} = 1} x_{in} \]
in the associated graded algebra of the May filtration
\[ E^0\A = \bigotimes_{i \ge 0,\, n \ge 1} \F[x_{in}]/(x_{in}^2). \]
For example $\Sq^3 = \Sq(3) \leftrightarrow x_{01}x_{11}$, $\Sq(0,1) \leftrightarrow x_{02}$ and
$\Sq(1,1) \leftrightarrow x_{01}x_{02}$. In $E^0\A$ a monomial is just a set of positions, so
divisibility, quotients and least common multiples are inclusion, difference and union of sets.
The positions relevant here, in increasing degree, are
\[ x_{01}\ (1),\quad x_{11}\ (2),\quad x_{02}\ (3),\quad x_{21}\ (4),\quad x_{12}\ (6),\quad x_{03}\ (7). \]

\paragraph{The monomial order.}
A term of a free module is $\Sq(R)\,g$ with $g$ a generator. Terms are compared by
\begin{enumerate}
\item the generator: the generator added later (of higher degree) comes first;
\item the May weight $\sum_{d_{i,n} = 1} (2n - 1)$: smaller comes first;
\item the digits read as a binary number, with the position of lowest degree most significant:
  smaller comes first.
\end{enumerate}
The \emph{leading term} $\LT(x)$ of an element is its first term in this order. Elements are
always printed with their leading term first. The point of this order is that the leading term
of a product $\Sq(R)\cdot x$ is $x^R \cdot \LT(x)$ computed in $E^0\A$, \emph{whenever that
product is non-zero}. Products of leading terms vanish exactly when a position occurs twice, and
those are the cases the algorithm treats separately (``singles'' below).

\paragraph{Names.}
The generator of $C_0$ is $\iota$. Generators of $C_s$, $s \ge 1$, are named after the classes of
$\Ext_\A^{s,t}(\F, \F)$ they represent: $[h_0], [h_1], [h_2] \in C_1$, $[h_0^2], [h_1^2],
[h_0h_2] \in C_2$, and so on. These names are only labels; the program numbers the generators of
each $C_s$ in the order they are found.

\section{The algorithm}

\paragraph{Levels.}
Level $s$ is a Gröbner basis $b^s_0, b^s_1, \dots$ of the graph of the differential,
\[ \Gamma_s = \{ (d y, y) : y \in C_{s+1} \} \subseteq C_s \oplus C_{s+1}, \]
for the order on the first component. Each $b^s_k = (x_1 \mid x_2)$ satisfies $d x_2 = x_1$,
and its leading term is that of $x_1$. The first components form a Gröbner basis of
$\operatorname{im}(d \colon C_{s+1} \to C_s)$, and pairs whose first component reduces to zero
produce elements of the kernel. The module being resolved enters as ``level $-1$'': the relations
$\Sq^{2^i}\iota = 0$ presenting $\F = \A/\A\{\Sq^1, \Sq^2, \Sq^4, \dots\}$.

\paragraph{Critical pairs and singles.}
When an element with leading term $L = x^{R}\,g$ joins level $s$, the program forms
\begin{itemize}
\item a \emph{critical pair} with each earlier element $b$ whose leading term $L'$ lies on the
  same generator: with $M = \operatorname{lcm}(L, L')$ in $E^0\A$, the S-element is
  $\Sq(M/L')\cdot b + \Sq(M/L)\cdot b_{\text{new}}$, in which the leading terms cancel. The pairs
  are pruned by the Gebauer--Möller criteria (divisibility and the chain criterion);
\item a \emph{critical single} $\Sq(x_{in}) \cdot b_{\text{new}}$ for each position $x_{in}$
  dividing $L$: then $x_{in} \cdot L = 0$ in $E^0\A$, so the leading term of the product is not
  predicted by the leading terms and the product must be reduced.
\end{itemize}
A pair or single has the degree of its leading monomial, and it is filed under that degree.

\paragraph{One degree.}
The program runs through $t = 1, 2, \dots$. In degree $t$:
\begin{enumerate}
\item Every level hands out its pairs and singles of degree $t$, and the relations of degree $t$
  are taken up.
\item Each of them is reduced: the first component modulo level $s$, carrying the second component
  along so that $d x_2 = x_1$ is preserved; then the second component modulo the first components
  of level $s+1$, i.e.\ modulo the boundaries $d(C_{s+2})$ known so far, which preserves
  $d x_2 = x_1$ because $dd = 0$. If the first component has become zero, the second component
  $z \in C_{s+1}$ is a cycle.
\item Going down from the top level, the results are compared with the elements found earlier in
  the same degree (``triangulation''), and then sorted:
  \begin{itemize}
  \item first component non-zero: a \textbf{new basis element} of level $s$;
  \item both components zero: \textbf{useless};
  \item first component zero, second component $z \ne 0$: $z$ is a cycle that is not a boundary
    (it is reduced modulo the whole image in degree $t$), so it becomes $d$ of a
    \textbf{new generator} $g$ of $C_{s+2}$, and $(z \mid g)$ joins level $s+1$. If $s + 2$
    exceeds the range, the outcome is \textbf{beyond} instead.
  \end{itemize}
\end{enumerate}
Because the levels are handled from the top down, the new elements of level $s+1$ in degree $t$
are known when the cycles of level $s$ are tested, which is what makes the resolution minimal.
No linear algebra is ever done: everything is rewriting of sparse sums of terms.

\paragraph{Truncation.}
A cycle found at level $s$ in degree $t$ is a generator of $C_{s+2}$ in stem $t-s-2$. Since we
want stems $\le 4$, level $s$ is only needed for $t \le s + 6$, and pairs of higher degree are
never formed.

\paragraph{Running it.}\sloppy
The trace was produced by a small driver around the public API of the crate
(\texttt{Resolution::sphere} with \texttt{Bounds \{ s\_max: 3, n\_max: 4 \}} at $p = 2$,
\texttt{SteenrodKind::Classical}), which calls \texttt{compute\_through(t)} for $t = 1, \dots, 7$
and records every reduction through a \texttt{PairObserver}.

\section{Step by step}

\subsection*{Degree $t = 1$}
\noindent The work of this degree: the relation $\Sq^{1}\iota$. All of them are reduced independently (in parallel in the program), and the results are then sorted from the top level down.
\medskip
\paragraph{Relation $\Sq^{1}\iota = 0$.} The module relation enters as the element $(0 \mid \Sq^{1}\iota)$ of level $-1$, so its $C_0$ component is a candidate for $d$ of a new generator. It is reduced modulo the image $\operatorname{im}(d\colon C_1\to C_0)$, i.e.\ the first components of level $0$.
No term of it is divisible by a leading term of level $0$ from earlier degrees.
\textbf{Outcome: new generator.} The relation $\Sq^{1}\iota$ is not a boundary, so $C_{1}$ gets a generator $[h_0]$ in degree $1$ with $d[h_0] = \Sq^{1}\iota$.
Level $0$ gets the element $b^{0}_{0} = \bigl(\Sq^{1}\iota \;\big|\; [h_0]\bigr)$, with leading term $x_{01}\,\iota$.
\medskip\noindent\emph{End of degree 1.} New basis elements: level $0$: $b^{0}_{0}$. New generators: $[h_0]\in C_{1}$.
\subsection*{Degree $t = 2$}
\noindent The work of this degree: the relation $\Sq^{2}\iota$; level $0$: single. All of them are reduced independently (in parallel in the program), and the results are then sorted from the top level down.
\medskip
\paragraph{Level $0$, single $x_{01}\cdot b^{0}_{0}$.} Since $x_{01}$ divides $\LT(b^{0}_{0}) = x_{01}\,\iota$ and $x_{in}^2 = 0$ in $E^0\mathcal A$, the leading term of $\Sq^{1}\cdot b^{0}_{0}$ is not the product of the leading terms, and the product has to be reduced:
\[\bigl(0 \;\big|\; \Sq^{1}[h_0]\bigr).\]
The first component is zero, so the second component $z = \Sq^{1}[h_0]$ is a cycle in $C_{1}$.
No term of it is divisible by a leading term of level $1$ from earlier degrees.
The cycle $\Sq^{1}[h_0]$ is set aside: cycles are tested once all results of level $0$ in this degree are sorted.
\medskip
\paragraph{The cycles of level $0$.} There is one.
\textbf{Outcome: new generator.} The cycle $\Sq^{1}[h_0]$ is not a boundary, so $C_{2}$ gets a generator $[h_0^2]$ in degree $2$ with $d[h_0^2] = \Sq^{1}[h_0]$.
Level $1$ gets the element $b^{1}_{0} = \bigl(\Sq^{1}[h_0] \;\big|\; [h_0^2]\bigr)$, with leading term $x_{01}\,[h_0]$.
\medskip
\paragraph{Relation $\Sq^{2}\iota = 0$.} The module relation enters as the element $(0 \mid \Sq^{2}\iota)$ of level $-1$, so its $C_0$ component is a candidate for $d$ of a new generator. It is reduced modulo the image $\operatorname{im}(d\colon C_1\to C_0)$, i.e.\ the first components of level $0$.
No term of it is divisible by a leading term of level $0$ from earlier degrees.
\textbf{Outcome: new generator.} The relation $\Sq^{2}\iota$ is not a boundary, so $C_{1}$ gets a generator $[h_1]$ in degree $2$ with $d[h_1] = \Sq^{2}\iota$.
Level $0$ gets the element $b^{0}_{1} = \bigl(\Sq^{2}\iota \;\big|\; [h_1]\bigr)$, with leading term $x_{11}\,\iota$.
\medskip\noindent\emph{End of degree 2.} New basis elements: level $0$: $b^{0}_{1}$; level $1$: $b^{1}_{0}$. New generators: $[h_1]\in C_{1}$, $[h_0^2]\in C_{2}$.
\subsection*{Degree $t = 3$}
\noindent The work of this degree: level $0$: pair; level $1$: single. All of them are reduced independently (in parallel in the program), and the results are then sorted from the top level down.
\medskip
\paragraph{Level $1$, single $x_{01}\cdot b^{1}_{0}$.} Since $x_{01}$ divides $\LT(b^{1}_{0}) = x_{01}\,[h_0]$ and $x_{in}^2 = 0$ in $E^0\mathcal A$, the leading term of $\Sq^{1}\cdot b^{1}_{0}$ is not the product of the leading terms, and the product has to be reduced:
\[\bigl(0 \;\big|\; \Sq^{1}[h_0^2]\bigr).\]
The first component is zero, so the second component $z = \Sq^{1}[h_0^2]$ is a cycle in $C_{2}$.
No term of it is divisible by a leading term of level $2$ from earlier degrees.
The cycle $\Sq^{1}[h_0^2]$ is set aside: cycles are tested once all results of level $1$ in this degree are sorted.
\medskip
\paragraph{The cycles of level $1$.} There is one.
\textbf{Outcome: new generator.} The cycle $\Sq^{1}[h_0^2]$ is not a boundary, so $C_{3}$ gets a generator $[h_0^3]$ in degree $3$ with $d[h_0^3] = \Sq^{1}[h_0^2]$.
Level $2$ gets the element $b^{2}_{0} = \bigl(\Sq^{1}[h_0^2] \;\big|\; [h_0^3]\bigr)$, with leading term $x_{01}\,[h_0^2]$.
\medskip
\paragraph{Level $0$, pair $(b^{0}_{0}, b^{0}_{1})$.} The leading terms $\LT(b^{0}_{0}) = x_{01}\,\iota$ and $\LT(b^{0}_{1}) = x_{11}\,\iota$ have least common multiple $x_{01}x_{11}\,\iota$, with cofactors $x_{11}$ and $x_{01}$. The S-element $\Sq^{2}\cdot b^{0}_{0} + \Sq^{1}\cdot b^{0}_{1}$, in which the leading terms cancel, is
\[\bigl(\Sq(0,1)\iota \;\big|\; \Sq^{1}[h_1] + \Sq^{2}[h_0]\bigr).\]
No term of the first component is divisible by a leading term of level $0$.
\textbf{Outcome: new basis element} $b^{0}_{2} = \bigl(\Sq(0,1)\iota \;\big|\; \Sq^{1}[h_1] + \Sq^{2}[h_0]\bigr)$ of level $0$, with leading term $x_{02}\,\iota$.
\medskip\noindent\emph{End of degree 3.} New basis elements: level $0$: $b^{0}_{2}$; level $2$: $b^{2}_{0}$. New generators: $[h_0^3]\in C_{3}$.
\subsection*{Degree $t = 4$}
\noindent The work of this degree: the relation $\Sq^{4}\iota$; level $0$: single, pair; level $2$: single. All of them are reduced independently (in parallel in the program), and the results are then sorted from the top level down.
\medskip
\paragraph{Level $2$, single $x_{01}\cdot b^{2}_{0}$.} Since $x_{01}$ divides $\LT(b^{2}_{0}) = x_{01}\,[h_0^2]$ and $x_{in}^2 = 0$ in $E^0\mathcal A$, the leading term of $\Sq^{1}\cdot b^{2}_{0}$ is not the product of the leading terms, and the product has to be reduced:
\[\bigl(0 \;\big|\; \Sq^{1}[h_0^3]\bigr).\]
The first component is zero, so the second component $z = \Sq^{1}[h_0^3]$ is a cycle in $C_{3}$.
There is no level $3$ (we stop at $s = 3$).
The cycle $\Sq^{1}[h_0^3]$ is set aside: cycles are tested once all results of level $2$ in this degree are sorted.
\medskip
\paragraph{The cycles of level $2$.} There is one.
\textbf{Outcome: beyond.} The cycle $\Sq^{1}[h_0^3]$ would give a generator of $C_{4}$, beyond $s = 3$, so it is dropped.
\medskip
\paragraph{Level $0$, single $x_{11}\cdot b^{0}_{1}$.} Since $x_{11}$ divides $\LT(b^{0}_{1}) = x_{11}\,\iota$ and $x_{in}^2 = 0$ in $E^0\mathcal A$, the leading term of $\Sq^{2}\cdot b^{0}_{1}$ is not the product of the leading terms, and the product has to be reduced:
\[\bigl(\Sq(1,1)\iota \;\big|\; \Sq^{2}[h_1]\bigr).\]
Reduce the first component modulo level $0$:
\begin{itemize}
\item The term $\Sq(1,1)\iota$ has leading monomial $x_{01}x_{02}\,\iota$, divisible by $\LT(b^{0}_{0})$ with quotient $x_{02}$. Adding $\Sq(0,1)\cdot b^{0}_{0} = \bigl(\Sq(1,1)\iota \;\big|\; \Sq(0,1)[h_0]\bigr)$ leaves $\bigl(0 \;\big|\; \Sq^{2}[h_1] + \Sq(0,1)[h_0]\bigr)$.
\end{itemize}
The first component is zero, so the second component $z = \Sq^{2}[h_1] + \Sq(0,1)[h_0]$ is a cycle in $C_{1}$.
No term of it is divisible by a leading term of level $1$ from earlier degrees.
The cycle $\Sq^{2}[h_1] + \Sq(0,1)[h_0]$ is set aside: cycles are tested once all results of level $0$ in this degree are sorted.
\medskip
\paragraph{Level $0$, pair $(b^{0}_{0}, b^{0}_{2})$.} The leading terms $\LT(b^{0}_{0}) = x_{01}\,\iota$ and $\LT(b^{0}_{2}) = x_{02}\,\iota$ have least common multiple $x_{01}x_{02}\,\iota$, with cofactors $x_{02}$ and $x_{01}$. The S-element $\Sq(0,1)\cdot b^{0}_{0} + \Sq^{1}\cdot b^{0}_{2}$, in which the leading terms cancel, is
\[\bigl(0 \;\big|\; \Sq^{3}[h_0] + \Sq(0,1)[h_0]\bigr).\]
The first component is zero, so the second component $z = \Sq^{3}[h_0] + \Sq(0,1)[h_0]$ is a cycle in $C_{1}$.
Reduce it modulo $\operatorname{im}(d\colon C_{2}\to C_{1})$, the first components of level $1$:
\begin{itemize}
\item The term $\Sq^{3}[h_0]$ has leading monomial $x_{01}x_{11}\,[h_0]$, divisible by $\LT(b^{1}_{0})$ with quotient $x_{11}$. Adding $\Sq^{2}\cdot \pi_1(b^{1}_{0}) = \Sq^{3}[h_0] + \Sq(0,1)[h_0]$ leaves $0$.
\end{itemize}
\textbf{Outcome: useless.} Both components vanish, and the reduction contributed nothing.
\medskip
\paragraph{The cycles of level $0$.} There is one.
\textbf{Outcome: new generator.} The cycle $\Sq^{2}[h_1] + \Sq(0,1)[h_0]$ is not a boundary, so $C_{2}$ gets a generator $[h_1^2]$ in degree $4$ with $d[h_1^2] = \Sq^{2}[h_1] + \Sq(0,1)[h_0]$.
Level $1$ gets the element $b^{1}_{1} = \bigl(\Sq^{2}[h_1] + \Sq(0,1)[h_0] \;\big|\; [h_1^2]\bigr)$, with leading term $x_{11}\,[h_1]$.
\medskip
\paragraph{Relation $\Sq^{4}\iota = 0$.} The module relation enters as the element $(0 \mid \Sq^{4}\iota)$ of level $-1$, so its $C_0$ component is a candidate for $d$ of a new generator. It is reduced modulo the image $\operatorname{im}(d\colon C_1\to C_0)$, i.e.\ the first components of level $0$.
No term of it is divisible by a leading term of level $0$ from earlier degrees.
\textbf{Outcome: new generator.} The relation $\Sq^{4}\iota$ is not a boundary, so $C_{1}$ gets a generator $[h_2]$ in degree $4$ with $d[h_2] = \Sq^{4}\iota$.
Level $0$ gets the element $b^{0}_{3} = \bigl(\Sq^{4}\iota \;\big|\; [h_2]\bigr)$, with leading term $x_{21}\,\iota$.
\medskip\noindent\emph{End of degree 4.} New basis elements: level $0$: $b^{0}_{3}$; level $1$: $b^{1}_{1}$. New generators: $[h_2]\in C_{1}$, $[h_1^2]\in C_{2}$.
\subsection*{Degree $t = 5$}
\noindent The work of this degree: level $0$: pair, pair. All of them are reduced independently (in parallel in the program), and the results are then sorted from the top level down.
\medskip
\paragraph{Level $0$, pair $(b^{0}_{1}, b^{0}_{2})$.} The leading terms $\LT(b^{0}_{1}) = x_{11}\,\iota$ and $\LT(b^{0}_{2}) = x_{02}\,\iota$ have least common multiple $x_{11}x_{02}\,\iota$, with cofactors $x_{02}$ and $x_{11}$. The S-element $\Sq(0,1)\cdot b^{0}_{1} + \Sq^{2}\cdot b^{0}_{2}$, in which the leading terms cancel, is
\[\bigl(0 \;\big|\; \Sq^{3}[h_1] + \Sq(1,1)[h_0]\bigr).\]
The first component is zero, so the second component $z = \Sq^{3}[h_1] + \Sq(1,1)[h_0]$ is a cycle in $C_{1}$.
Reduce it modulo $\operatorname{im}(d\colon C_{2}\to C_{1})$, the first components of level $1$:
\begin{itemize}
\item The term $\Sq^{3}[h_1]$ has leading monomial $x_{01}x_{11}\,[h_1]$, divisible by $\LT(b^{1}_{1})$ with quotient $x_{01}$. Adding $\Sq^{1}\cdot \pi_1(b^{1}_{1}) = \Sq^{3}[h_1] + \Sq(1,1)[h_0]$ leaves $0$.
\end{itemize}
\textbf{Outcome: useless.} Both components vanish, and the reduction contributed nothing.
\medskip
\paragraph{Level $0$, pair $(b^{0}_{0}, b^{0}_{3})$.} The leading terms $\LT(b^{0}_{0}) = x_{01}\,\iota$ and $\LT(b^{0}_{3}) = x_{21}\,\iota$ have least common multiple $x_{01}x_{21}\,\iota$, with cofactors $x_{21}$ and $x_{01}$. The S-element $\Sq^{4}\cdot b^{0}_{0} + \Sq^{1}\cdot b^{0}_{3}$, in which the leading terms cancel, is
\[\bigl(\Sq(2,1)\iota \;\big|\; \Sq^{1}[h_2] + \Sq^{4}[h_0]\bigr).\]
Reduce the first component modulo level $0$:
\begin{itemize}
\item The term $\Sq(2,1)\iota$ has leading monomial $x_{11}x_{02}\,\iota$, divisible by $\LT(b^{0}_{1})$ with quotient $x_{02}$. Adding $\Sq(0,1)\cdot b^{0}_{1} = \bigl(\Sq(2,1)\iota \;\big|\; \Sq(0,1)[h_1]\bigr)$ leaves $\bigl(0 \;\big|\; \Sq^{1}[h_2] + \Sq(0,1)[h_1] + \Sq^{4}[h_0]\bigr)$.
\end{itemize}
The first component is zero, so the second component $z = \Sq^{1}[h_2] + \Sq(0,1)[h_1] + \Sq^{4}[h_0]$ is a cycle in $C_{1}$.
No term of it is divisible by a leading term of level $1$ from earlier degrees.
The cycle $\Sq^{1}[h_2] + \Sq(0,1)[h_1] + \Sq^{4}[h_0]$ is set aside: cycles are tested once all results of level $0$ in this degree are sorted.
\medskip
\paragraph{The cycles of level $0$.} There is one.
\textbf{Outcome: new generator.} The cycle $\Sq^{1}[h_2] + \Sq(0,1)[h_1] + \Sq^{4}[h_0]$ is not a boundary, so $C_{2}$ gets a generator $[h_0h_2]$ in degree $5$ with $d[h_0h_2] = \Sq^{1}[h_2] + \Sq(0,1)[h_1] + \Sq^{4}[h_0]$.
Level $1$ gets the element $b^{1}_{2} = \bigl(\Sq^{1}[h_2] + \Sq(0,1)[h_1] + \Sq^{4}[h_0] \;\big|\; [h_0h_2]\bigr)$, with leading term $x_{01}\,[h_2]$.
\medskip\noindent\emph{End of degree 5.} New basis elements: level $1$: $b^{1}_{2}$. New generators: $[h_0h_2]\in C_{2}$.
\subsection*{Degree $t = 6$}
\noindent The work of this degree: level $0$: single, pair; level $1$: single, single. All of them are reduced independently (in parallel in the program), and the results are then sorted from the top level down.
\medskip
\paragraph{Level $1$, single $x_{11}\cdot b^{1}_{1}$.} Since $x_{11}$ divides $\LT(b^{1}_{1}) = x_{11}\,[h_1]$ and $x_{in}^2 = 0$ in $E^0\mathcal A$, the leading term of $\Sq^{2}\cdot b^{1}_{1}$ is not the product of the leading terms, and the product has to be reduced:
\[\bigl(\Sq(1,1)[h_1] + \Sq(2,1)[h_0] \;\big|\; \Sq^{2}[h_1^2]\bigr).\]
No term of the first component is divisible by a leading term of level $1$.
\textbf{Outcome: new basis element} $b^{1}_{3} = \bigl(\Sq(1,1)[h_1] + \Sq(2,1)[h_0] \;\big|\; \Sq^{2}[h_1^2]\bigr)$ of level $1$, with leading term $x_{01}x_{02}\,[h_1]$.
\medskip
\paragraph{Level $1$, single $x_{01}\cdot b^{1}_{2}$.} Since $x_{01}$ divides $\LT(b^{1}_{2}) = x_{01}\,[h_2]$ and $x_{in}^2 = 0$ in $E^0\mathcal A$, the leading term of $\Sq^{1}\cdot b^{1}_{2}$ is not the product of the leading terms, and the product has to be reduced:
\[\bigl(\Sq(1,1)[h_1] + \Sq^{5}[h_0] \;\big|\; \Sq^{1}[h_0h_2]\bigr).\]
Reduce the first component modulo level $1$:
\begin{itemize}
\item The term $\Sq^{5}[h_0]$ has leading monomial $x_{01}x_{21}\,[h_0]$, divisible by $\LT(b^{1}_{0})$ with quotient $x_{21}$. Adding $\Sq^{4}\cdot b^{1}_{0} = \bigl(\Sq^{5}[h_0] + \Sq(2,1)[h_0] \;\big|\; \Sq^{4}[h_0^2]\bigr)$ leaves $\bigl(\Sq(1,1)[h_1] + \Sq(2,1)[h_0] \;\big|\; \Sq^{1}[h_0h_2] + \Sq^{4}[h_0^2]\bigr)$.
\end{itemize}
Its first component contains $\LT(b^{1}_{3})$ of the element $b^{1}_{3}$ just found in this degree; adding it gives $\bigl(0 \;\big|\; \Sq^{1}[h_0h_2] + \Sq^{2}[h_1^2] + \Sq^{4}[h_0^2]\bigr)$.
The cycle $\Sq^{1}[h_0h_2] + \Sq^{2}[h_1^2] + \Sq^{4}[h_0^2]$ is set aside: cycles are tested once all results of level $1$ in this degree are sorted.
\medskip
\paragraph{The cycles of level $1$.} There is one.
\textbf{Outcome: new generator.} The cycle $\Sq^{1}[h_0h_2] + \Sq^{2}[h_1^2] + \Sq^{4}[h_0^2]$ is not a boundary, so $C_{3}$ gets a generator $[h_0^2h_2]$ in degree $6$ with $d[h_0^2h_2] = \Sq^{1}[h_0h_2] + \Sq^{2}[h_1^2] + \Sq^{4}[h_0^2]$.
Level $2$ gets the element $b^{2}_{1} = \bigl(\Sq^{1}[h_0h_2] + \Sq^{2}[h_1^2] + \Sq^{4}[h_0^2] \;\big|\; [h_0^2h_2]\bigr)$, with leading term $x_{01}\,[h_0h_2]$.
\medskip
\paragraph{Level $0$, single $x_{02}\cdot b^{0}_{2}$.} Since $x_{02}$ divides $\LT(b^{0}_{2}) = x_{02}\,\iota$ and $x_{in}^2 = 0$ in $E^0\mathcal A$, the leading term of $\Sq(0,1)\cdot b^{0}_{2}$ is not the product of the leading terms, and the product has to be reduced:
\[\bigl(0 \;\big|\; \Sq(1,1)[h_1] + \Sq(2,1)[h_0]\bigr).\]
The first component is zero, so the second component $z = \Sq(1,1)[h_1] + \Sq(2,1)[h_0]$ is a cycle in $C_{1}$.
No term of it is divisible by a leading term of level $1$ from earlier degrees.
The second component contains the leading term of $b^{1}_{3}$, found in this degree; adding its first component leaves $0$.
\textbf{Outcome: useless.} The cycle is the first component of $b^{1}_{3}$, found one level up earlier in this degree, so it is the boundary $d\bigl(\Sq^{2}[h_1^2]\bigr)$. Both components now vanish, and the reduction contributed nothing.
\medskip
\paragraph{Level $0$, pair $(b^{0}_{1}, b^{0}_{3})$.} The leading terms $\LT(b^{0}_{1}) = x_{11}\,\iota$ and $\LT(b^{0}_{3}) = x_{21}\,\iota$ have least common multiple $x_{11}x_{21}\,\iota$, with cofactors $x_{21}$ and $x_{11}$. The S-element $\Sq^{4}\cdot b^{0}_{1} + \Sq^{2}\cdot b^{0}_{3}$, in which the leading terms cancel, is
\[\bigl(\Sq(0,2)\iota \;\big|\; \Sq^{2}[h_2] + \Sq^{4}[h_1]\bigr).\]
No term of the first component is divisible by a leading term of level $0$.
\textbf{Outcome: new basis element} $b^{0}_{4} = \bigl(\Sq(0,2)\iota \;\big|\; \Sq^{2}[h_2] + \Sq^{4}[h_1]\bigr)$ of level $0$, with leading term $x_{12}\,\iota$.
\medskip\noindent\emph{End of degree 6.} New basis elements: level $0$: $b^{0}_{4}$; level $1$: $b^{1}_{3}$; level $2$: $b^{2}_{1}$. New generators: $[h_0^2h_2]\in C_{3}$.
\subsection*{Degree $t = 7$}
\noindent Level $0$ is no longer needed: a cycle found at level $s$ in this degree would be a generator of $C_{s+2}$ in stem $7-s-2 > 4$, so its pairs are not formed.
\noindent The work of this degree: level $1$: single; level $2$: single. All of them are reduced independently (in parallel in the program), and the results are then sorted from the top level down.
\medskip
\paragraph{Level $2$, single $x_{01}\cdot b^{2}_{1}$.} Since $x_{01}$ divides $\LT(b^{2}_{1}) = x_{01}\,[h_0h_2]$ and $x_{in}^2 = 0$ in $E^0\mathcal A$, the leading term of $\Sq^{1}\cdot b^{2}_{1}$ is not the product of the leading terms, and the product has to be reduced:
\[\bigl(\Sq^{3}[h_1^2] + \Sq^{5}[h_0^2] \;\big|\; \Sq^{1}[h_0^2h_2]\bigr).\]
Reduce the first component modulo level $2$:
\begin{itemize}
\item The term $\Sq^{5}[h_0^2]$ has leading monomial $x_{01}x_{21}\,[h_0^2]$, divisible by $\LT(b^{2}_{0})$ with quotient $x_{21}$. Adding $\Sq^{4}\cdot b^{2}_{0} = \bigl(\Sq^{5}[h_0^2] + \Sq(2,1)[h_0^2] \;\big|\; \Sq^{4}[h_0^3]\bigr)$ leaves $\bigl(\Sq^{3}[h_1^2] + \Sq(2,1)[h_0^2] \;\big|\; \Sq^{1}[h_0^2h_2] + \Sq^{4}[h_0^3]\bigr)$.
\end{itemize}
\textbf{Outcome: new basis element} $b^{2}_{2} = \bigl(\Sq^{3}[h_1^2] + \Sq(2,1)[h_0^2] \;\big|\; \Sq^{1}[h_0^2h_2] + \Sq^{4}[h_0^3]\bigr)$ of level $2$, with leading term $x_{01}x_{11}\,[h_1^2]$.
\medskip
\paragraph{Level $1$, single $x_{01}\cdot b^{1}_{3}$.} Since $x_{01}$ divides $\LT(b^{1}_{3}) = x_{01}x_{02}\,[h_1]$ and $x_{in}^2 = 0$ in $E^0\mathcal A$, the leading term of $\Sq^{1}\cdot b^{1}_{3}$ is not the product of the leading terms, and the product has to be reduced:
\[\bigl(\Sq(3,1)[h_0] \;\big|\; \Sq^{3}[h_1^2]\bigr).\]
Reduce the first component modulo level $1$:
\begin{itemize}
\item The term $\Sq(3,1)[h_0]$ has leading monomial $x_{01}x_{11}x_{02}\,[h_0]$, divisible by $\LT(b^{1}_{0})$ with quotient $x_{11}x_{02}$. Adding $\Sq(2,1)\cdot b^{1}_{0} = \bigl(\Sq(3,1)[h_0] \;\big|\; \Sq(2,1)[h_0^2]\bigr)$ leaves $\bigl(0 \;\big|\; \Sq^{3}[h_1^2] + \Sq(2,1)[h_0^2]\bigr)$.
\end{itemize}
The first component is zero, so the second component $z = \Sq^{3}[h_1^2] + \Sq(2,1)[h_0^2]$ is a cycle in $C_{2}$.
No term of it is divisible by a leading term of level $2$ from earlier degrees.
The second component contains the leading term of $b^{2}_{2}$, found in this degree; adding its first component leaves $0$.
\textbf{Outcome: useless.} The cycle is the first component of $b^{2}_{2}$, found one level up earlier in this degree, so it is the boundary $d\bigl(\Sq^{1}[h_0^2h_2] + \Sq^{4}[h_0^3]\bigr)$. Both components now vanish, and the reduction contributed nothing.
\medskip\noindent\emph{End of degree 7.} New basis elements: level $2$: $b^{2}_{2}$. No new generators.

\section{The result}

\paragraph{Generators and differentials.}
\begin{center}
\begin{tabular}{lcccl}
\toprule
generator & $s$ & $t$ & $t-s$ & differential \\
\midrule
$\iota$ & 0 & 0 & 0 & --- \\
$[h_0]$ & 1 & 1 & 0 & $\Sq^{1}\iota$ \\
$[h_1]$ & 1 & 2 & 1 & $\Sq^{2}\iota$ \\
$[h_0^2]$ & 2 & 2 & 0 & $\Sq^{1}[h_0]$ \\
$[h_0^3]$ & 3 & 3 & 0 & $\Sq^{1}[h_0^2]$ \\
$[h_2]$ & 1 & 4 & 3 & $\Sq^{4}\iota$ \\
$[h_1^2]$ & 2 & 4 & 2 & $\Sq^{2}[h_1] + \Sq(0,1)[h_0]$ \\
$[h_0h_2]$ & 2 & 5 & 3 & $\Sq^{1}[h_2] + \Sq(0,1)[h_1] + \Sq^{4}[h_0]$ \\
$[h_0^2h_2]$ & 3 & 6 & 3 & $\Sq^{1}[h_0h_2] + \Sq^{2}[h_1^2] + \Sq^{4}[h_0^2]$ \\
\bottomrule
\end{tabular}
\end{center}
Since the resolution is minimal, the generators are a basis of $\Ext_\A^{s,t}(\F,\F)$ in this
range:
\begin{center}
\begin{tabular}{c|ccccc}
$s \backslash t-s$ & 0 & 1 & 2 & 3 & 4 \\
\midrule
$3$ & $h_0^3$ &  &  & $h_0^2h_2$ &  \\
$2$ & $h_0^2$ &  & $h_1^2$ & $h_0h_2$ &  \\
$1$ & $h_0$ & $h_1$ &  & $h_2$ &  \\
$0$ & $1$ &  &  &  &  \\
\end{tabular}
\end{center}
This is the familiar picture: the $h_0$-tower in stem $0$, $h_1$ and $h_1^2$ in stems $1$ and
$2$, and $h_2$, $h_0h_2$, $h_0^2h_2 = h_1^3$ in stem $3$; stem $4$ is empty.
The cycle beyond $s = 3$ found in degree $4$ is the next class $h_0^4$ of the tower.

\paragraph{The Gröbner bases.}
\begin{center}
\small
\begin{longtable}{lclll}
\toprule
element & $t$ & leading term & first component $x_1$ & second component $x_2$ \\
\midrule
$b^{0}_{0}$ & 1 & $x_{01}\,\iota$ & $\Sq^{1}\iota$ & $[h_0]$ \\
$b^{0}_{1}$ & 2 & $x_{11}\,\iota$ & $\Sq^{2}\iota$ & $[h_1]$ \\
$b^{0}_{2}$ & 3 & $x_{02}\,\iota$ & $\Sq(0,1)\iota$ & $\Sq^{1}[h_1] + \Sq^{2}[h_0]$ \\
$b^{0}_{3}$ & 4 & $x_{21}\,\iota$ & $\Sq^{4}\iota$ & $[h_2]$ \\
$b^{0}_{4}$ & 6 & $x_{12}\,\iota$ & $\Sq(0,2)\iota$ & $\Sq^{2}[h_2] + \Sq^{4}[h_1]$ \\
$b^{1}_{0}$ & 2 & $x_{01}\,[h_0]$ & $\Sq^{1}[h_0]$ & $[h_0^2]$ \\
$b^{1}_{1}$ & 4 & $x_{11}\,[h_1]$ & $\Sq^{2}[h_1] + \Sq(0,1)[h_0]$ & $[h_1^2]$ \\
$b^{1}_{2}$ & 5 & $x_{01}\,[h_2]$ & $\Sq^{1}[h_2] + \Sq(0,1)[h_1] + \Sq^{4}[h_0]$ & $[h_0h_2]$ \\
$b^{1}_{3}$ & 6 & $x_{01}x_{02}\,[h_1]$ & $\Sq(1,1)[h_1] + \Sq(2,1)[h_0]$ & $\Sq^{2}[h_1^2]$ \\
$b^{2}_{0}$ & 3 & $x_{01}\,[h_0^2]$ & $\Sq^{1}[h_0^2]$ & $[h_0^3]$ \\
$b^{2}_{1}$ & 6 & $x_{01}\,[h_0h_2]$ & $\Sq^{1}[h_0h_2] + \Sq^{2}[h_1^2] + \Sq^{4}[h_0^2]$ & $[h_0^2h_2]$ \\
$b^{2}_{2}$ & 7 & $x_{01}x_{11}\,[h_1^2]$ & $\Sq^{3}[h_1^2] + \Sq(2,1)[h_0^2]$ & $\Sq^{1}[h_0^2h_2] + \Sq^{4}[h_0^3]$ \\
\bottomrule
\end{longtable}
\end{center}

\paragraph{Counters.}
The statistics the program keeps per level ($s = -1$ are the module relations):
\begin{center}
\footnotesize\setlength{\tabcolsep}{3.5pt}
\begin{tabular}{rrrrrrrrrrrr}
\toprule
$s$ & formed & divisible & chain & singles & reduced & useless & elements & boundary & gens & beyond & steps \\
\midrule
-1 & 0 & 0 & 0 & 0 & 3 & 0 & 0 & 0 & 3 & 0 & 0 \\
0 & 5 & 0 & 0 & 3 & 8 & 3 & 2 & 0 & 3 & 0 & 4 \\
1 & 0 & 0 & 0 & 4 & 4 & 1 & 1 & 0 & 2 & 0 & 2 \\
2 & 0 & 0 & 0 & 2 & 2 & 0 & 1 & 0 & 0 & 1 & 1 \\
all & 5 & 0 & 0 & 9 & 17 & 4 & 4 & 0 & 8 & 1 & 7 \\
\bottomrule
\end{tabular}
\end{center}
Here ``formed'' counts critical pairs, ``divisible'' and ``chain'' those removed by the
Gebauer--Möller criteria (none in this small range), and ``steps'' the reduction steps. Of the
17 reductions, 4 were useless and one was beyond the range; the other 12 each produced a basis
element or a generator.

\end{document}
